\documentclass{article}
\usepackage[utf8]{inputenc}
\usepackage[T1]{fontenc}
\usepackage{geometry}
\geometry{
    a4paper,
    left=2.5cm,
    right=2.5cm,
    top=2.5cm,
    bottom=2.5cm
}
\usepackage{graphicx}
\usepackage{float}
\usepackage{hyperref}
\usepackage[english]{babel}
\usepackage{amsmath}
\usepackage{amssymb}
\usepackage{array}
\usepackage{booktabs}
\usepackage{caption}
\usepackage{subcaption}
\usepackage{enumitem}
\usepackage{fancyhdr}
\usepackage{titlesec}
\usepackage{url}
\usepackage{textcomp}
\usepackage{xcolor}
\usepackage{listings}
\usepackage{multirow}
\usepackage{longtable}
\usepackage{tabularx}
\usepackage{times}

% Set up page style
\pagestyle{fancy}
\fancyhf{}
\rfoot{\thepage}
\renewcommand{\headrulewidth}{0pt}

% Title formatting
\titleformat{\chapter}[display]{\normalfont\bfseries}{\chaptertitle\thechapter}{0pt}{\Huge}
\titlespacing*{\chapter}{0pt}{50pt}{40pt}

% Document info
\title{Template Document}
\author{Academic Support Agent}
\date{\today}

\begin{document}
\maketitle

{
\setbeamercolor{background canvas}{bg=NavyBlue!50!DarkOliveGreen, fg=white}
\setbeamercolor{normal text}{fg=white}
\maketitle
}%This is the colour of the first slide. bg= background and fg=foreground

\metroset{titleformat frame=smallcaps} %This changes the titles for small caps

\begin{frame}{Outline}
  \setbeamertemplate{section in toc}[sections numbered] %This is numbering the sections
  
\end{frame}

\begin{frame}{Objectives}
\underline{\textsc{Some text:}}
\begin{small}
This is some small Text. 
\end{small}

\metroset{block=fill}
\begin{exampleblock}{\textsc{Example block}}
\begin{itemize}
    \item You know how to do itemize
    \item Also here
\end{itemize}
\end{exampleblock}
\end{frame}


\section{Introduction}

\begin{frame}[fragile]{Introduction: blah blah} %You can change fragile by standout
Text Text Text Text. \\You can change the size of the footnote text like  Text\footnote{\small{ here.}} Text\footnote{\large{And here.}} Text\footnote{\tiny{And here.}}
\begin{itemize} %The symbol of the items can be changed by which ever you want, this is just an example.
    \item[$\diamond$] Text,
    \item[$\diamond$] Text,
    \item[$\diamond$] Text.
\end{itemize}
An equation without number could be represented by:
\begin{equation*}
    c^{2} = a^{2} + b^{2}
\end{equation*}
That's all for this slide.
\end{frame}

\begin{frame}[standout]{This is other type of slide}
There is some text here.
And an equation with number:
\begin{equation}
    E^{2} = m^{2} + p^{2}
\end{equation}
\end{frame}

% Content from Othersection would be included here


\section{This is another section}
\begin{frame}{Frame Title} %You can also not write fragile or standout and you can see how it looks
    Hello world!
\end{frame}

\section{Final section}

\begin{frame}{Conclusion}
    These are the final words, you do your best to try to wake up everyone that was listening to your talk.
\end{frame}

{\setbeamercolor{palette primary}{fg=black, bg=orange!30} %You can change the colours
\begin{frame}[standout]
  Thank you! And thank to yourself because you did all the job. 
\end{frame}
}


\begin{frame}{Back up}
    These slides won't appear in the table of contents and will not be counted as the total slides.
\end{frame}
\end{document}